\section{Methodology}
\label{sec:method}

The primary objective of this section is to provide a complete, self-contained specification enabling independent reproduction of our experimental workflow and evaluation.

\subsection{Dataset}
\label{subsec:dataset}
We evaluate our system on the \textit{TestExplora} benchmark~\cite{testexplora}, a state-of-the-art dataset for documentation-derived test generation available at \url{https://github.com/microsoft/TestExplora}. The full benchmark comprises $N=1,552$ unique tasks across 482 public Python repositories created post-2023 to prevent data contamination in pre-trained LLMs. Each task provides a defective codebase snapshot paired with documentation-derived intent validated via the automated \textit{DocAgent} oracle. 

To evaluate the complete dataset without domain sampling bias, we included all $N=1,552$ tasks. The tasks span 10 software domain categories, including General Software Utilities (1,281 tasks), Database Engines (47 tasks), Infrastructure Monitoring (47 tasks), Web Systems (42 tasks), Data Science/ML (38 tasks), and Linters (30 tasks).

\subsection{Experimental Setup and Pipeline}
\label{subsec:setup}
Our system implements a Two-Agent Paradigm using OpenRouter API endpoints:
\begin{enumerate}
    \item \textbf{Exploration Agent:} Powered by \texttt{DeepSeek-V3} (\texttt{deepseek/deepseek-chat}) with $\text{temperature}=0.0$ and $\text{max\_tokens}=4096$, operating via the SWE-agent CLI bash terminal interface for up to 80 interaction turns per task.
    \item \textbf{Code Action Fixer Agent:} Powered by \texttt{Llama-3.3-70B-Instruct} (\texttt{meta-llama/llama-3.3-70b-instruct}) with $\text{temperature}=0.0$ and $\text{max\_tokens}=4096$, performing logical deduction to clean terminal logs and synthesize standalone Python test patches.
\end{enumerate}

The System Prompts for both agents are specified below:

\begin{lstlisting}[caption={System Prompt for Exploration Agent}, label={lst:prompt_explore}]
You are an expert autonomous software testing agent equipped with an interactive bash terminal.
Your task is to proactively discover latent bugs in the given Python repository based on documentation intent.
Instructions:
1. Explore the codebase using directory navigation and file viewing commands.
2. Draft candidate unit tests and execute pytest to capture runtime execution logs.
3. Observe tracebacks, identify compliance biases, and write adversarial test inputs.
4. Continue iterating until you isolate clear failing behavior on the buggy implementation.
\end{lstlisting}

\begin{lstlisting}[caption={System Prompt for Code Action Fixer Agent}, label={lst:prompt_fixer}]
Below is the execution log and candidate test draft produced by the exploration agent.
Analyze the execution log carefully. Your task is to refine and generate a clean, standalone Python test suite patch.
Requirements:
1. Ensure assertions directly capture the documentation intent.
2. Fix syntax errors and invalid test setups.
3. Output ONLY valid executable Python code for the final test suite.
\end{lstlisting}

\textbf{Execution Pipeline:} Input task $\rightarrow$ 16-Shard Round-Robin Distribution ($\text{Task}_i \in \text{Shard}_k \iff i \pmod{16} = k$) $\rightarrow$ Exploration Agent SWE-agent terminal execution (80 turns) $\rightarrow$ Execution log extraction $\rightarrow$ Code Action Fixer synthesis $\rightarrow$ DocAgent Oracle verdict evaluation.

\textbf{Edge Case, Cost, and Fault Handling:} To handle API rate limits (HTTP 429), server overloads (HTTP 503), and connection dropouts, we engineered an exponential backoff retry mechanism with randomized jitter: $t_{\text{wait}} = 2^{\text{attempt}} + \text{Uniform}(0, 1)$ seconds (up to 5 retries). Shell command executions were constrained by a 30-second hard timeout. The total execution cost for all 4,656 agent runs across $N=1,552$ tasks ($K=3$ runs) was approximately \$10--\$12 USD via paid OpenRouter API endpoints.

Our replication package, including datasets, prompts, scripts, and raw results, is publicly available at: \url{https://github.com/lctx9/group4}.

\subsection{Metrics}
\label{subsec:metrics}
We evaluate test suite quality using the Fail-to-Pass (F2P) rate metric defined by TestExplora~\cite{testexplora}:
\begin{equation}
\text{F2P Rate} = \frac{\sum_{i=1}^{N} \mathbb{I}(\text{Test}_i(\text{Code}_{\text{buggy}}) = \text{Fail} \land \text{Test}_i(\text{Code}_{\text{fixed}}) = \text{Pass})}{N} \times 100
\end{equation}
A test suite achieves F2P success if and only if it fails on defective code (proving bug detection) and passes on corrected code (proving assertion validity). Single-run performance is reported as Pass@1. Multi-attempt cumulative success across $K=3$ independent sampling runs is evaluated as Pass@3:
\begin{equation}
\text{Pass@}k = \frac{\sum_{i=1}^{N} \mathbb{I}\left( \exists j \in \{1,\dots,k\}: \text{F2P}_{i,j} = \text{Success} \right)}{N} \times 100
\end{equation}
The static Plain Context baseline threshold is set to 16.06\%, established by TestExplora~\cite{testexplora}.

\subsection{Statistical Analysis Plan}
\label{subsec:stat}
Following guidelines for statistical testing in software engineering~\cite{arcuri2011practical}, we employ non-parametric tests due to the non-normal distribution of repository F2P rates:
\begin{itemize}
    \item \textbf{Statistical Hypothesis Test:} Paired one-tailed Wilcoxon signed-rank test ($W$) evaluated at the repository level ($N_{\text{obs}} = 1,311$) using \texttt{scipy.stats.wilcoxon} with significance threshold $\alpha = 0.05$.
    \item \textbf{Effect Size Measurement:} Non-parametric Cliff's Delta ($\delta$) calculated over paired repository observations. Effect size magnitudes are interpreted according to standard thresholds~\cite{arcuri2011practical, cliffs}: Negligible ($|\delta| < 0.147$), Small ($0.147 \le |\delta| < 0.33$), Medium ($0.33 \le |\delta| < 0.474$), and Large ($|\delta| \ge 0.474$).
\end{itemize}
